\section{How much is attainable, and where it comes from}
\label{sec:benchmark}

This section answers RQ1 and RQ2. Nothing in it is propagated: every quantity
is computed from the tabulated defect and the reference state-transition
matrix, integrated with the variational equations along the reference arc,
so the comparison is free of integrator noise and of the step-size effects the
previous campaign had to control. What it therefore establishes is a
first-order statement about a fixed reference trajectory;
Section~\ref{sec:campaign} tests whether the propagated arcs agree.

\paragraph{Predicted error, not measured error.} Every error in this section is
the first-order prediction $\hat E=\sqrt{J}$ of Eq.~\eqref{eq:quadratic-form},
and every ratio is written $\hat\rho$. The propagated quantities $E$ and $\rho$
appear only from Section~\ref{sec:campaign} onward. The two are kept
typographically distinct throughout because they answer different questions and
are produced at different stages: no hypothesis stated on $\hat E$ is reported
as though it had been tested on $E$. Where the resolution rule is invoked here
it uses each orbit's archived numerical envelope, which is external to this
paper and pre-existing, since no propagation of this campaign has yet occurred
to supply one.

\subsection{The two separable allocations}
\label{sec:benchmark-sens}

Two of the four rungs below are separable and are stated here once, because
the Frank--Wolfe subproblem of Section~\ref{sec:algorithm-certificate} has the
same structure. Both are per-interval minimizations with a single multiplier
$\lambda$ bisected to the budget:

\begin{equation}
  \Aforce:\;\;
  \argmin_{N\in\mathcal{N}}\Bigl[\textstyle\sum_{i\in I_q}
    \lVert\Delta\mathbf a(t_i,N)\rVert^2\Delta t_i
    +\lambda W_q N^2\Bigr],
  \qquad
  \Asens:\;\;
  \argmin_{N\in\mathcal{N}}\Bigl[\textstyle\sum_{i\in I_q}
    \Delta t_i\,\Delta\mathbf a(t_i,N)^{\top}\mathbf K_i\,
    \Delta\mathbf a(t_i,N)
    +\lambda W_q N^2\Bigr].
  \label{eq:separable}
\end{equation}

\Aforce\ is the force-level benchmark of the previous paper: it never sees
$\Phi$ at all. \Asens\ is the natural sensitivity-weighted repair, and it is
worth being precise about where its weight comes from, because two wrong
answers are close at hand. Borrowing a scalar such as
$\lVert\Phi(T,t)\mathbf B\rVert^2$ from the terminal-state proxy would be an
arbitrary import, and matched to the wrong metric besides. Taking the discrete
diagonal $\mathbf u_i^{\top}\mathbf Q_{ii}\mathbf u_i$ literally would make the
criterion depend on the accumulation resolution, for the reason given in
Section~\ref{sec:problem-discretization}. The weight is instead the local
sensitivity kernel $\mathbf K_i$ of Eq.~\eqref{eq:local-sensitivity} --- the
objective's own coupling evaluated on its diagonal, grid-invariant, and matched
to the arc-RMS metric rather than to a terminal one.

The ladder is therefore a decomposition of one trajectory kernel and not a
comparison of three unrelated criteria: \Aforce\ ignores it, \Asens\ retains
its local diagonal $\mathbf K_i$, \Asign\ retains the full cross-epoch
coupling. RQ2 reduces to a single question --- how
much of the attainable gain lives off the diagonal --- and $G_{\mathrm{sens}}$
and $G_{\mathrm{sign}}$ below measure the two steps of that decomposition.

\subsection{The ladder}

Four allocations are compared at the same budget, each carrying one more piece
of the dynamics than the last.

\begin{table}[!htbp]
  \centering\small
  \caption{Allocation ladder at $\beta=1$. Entries are
  $\hat\rho(\Fop,\text{policy})=\hat E_{\Fop}/\hat E_{\text{policy}}$ on the
  first-order prediction, medians of per-orbit ratios, so values above one
  favour the policy. All four are evaluated on the same reference trajectories
  and the same tabulated defect, and none is propagated.
  \ph{fill from WP2, WP3}}
  \label{tab:ladder}
  \begin{tabular}{@{}llcccc@{}}
    \toprule
    & Allocation & Information used & A & B & C \\
    \midrule
    1 & \Rint\  & radius & \ph{value} & \ph{value} & \ph{value} \\
    2 & \Aforce & + reference field, pointwise & \ph{value} & \ph{value} & \ph{value} \\
    3 & \Asens  & + diagonal kernel $\mathbf K_i$ & \ph{value} & \ph{value} & \ph{value} \\
    4 & \Asign  & + signed coupling $\mathbf Q$ & \ph{value} & \ph{value} & \ph{value} \\
    \bottomrule
  \end{tabular}
\end{table}

\paragraph{RQ1 --- the attainable gain.} \ph{one paragraph} \Asign\ sits a
median \ph{value} above \Fop\ and \ph{value} above \Rint, the strongest
deployable schedule previously available. \ph{state whether the gain clears the
numerical envelope on how many orbits, and name the orbits where it does not.}

\dnote{Gate G1 of \code{../ROADMAP.md} is decided here. If
$\hat\rho(\Rint,\Asign)<1.5$ at the design median, this section becomes the
paper's principal result in the negative direction and
Sections~\ref{sec:controller} and~\ref{sec:campaign} contract sharply.}

\paragraph{RQ2 --- weight against sign.} The two increments are reported
separately because they are different claims. Raw error differences are not
comparable across orbits of different scale, so each increment is a log ratio:

\begin{equation}
  G_{\mathrm{sens}} = \log\frac{\hat E_{\Aforce}}{\hat E_{\Asens}},
  \qquad
  G_{\mathrm{sign}} = \log\frac{\hat E_{\Asens}}{\hat E_{\Asign}} ,
  \label{eq:increments}
\end{equation}

the first being what correcting the weighting buys and the second what the
coupling buys on top of it. \ph{one paragraph with the two medians and their
ratio.} The previous paper predicted that reweighting alone would not close the
gap; \ph{state whether that is confirmed, and by how much.}

\begin{figure}[!tbp]
  \centering
  \ph{figures/fig\_oa\_ladder.pdf}
  \caption{The allocation ladder, one point per orbit, designs pooled. Bars are
  medians. Rungs 1--2 use magnitudes only; rung 3 adds the sensitivity weight;
  rung 4 adds the signed coupling.}
  \label{fig:ladder}
\end{figure}

\subsection{What the schedules look like}

\ph{one paragraph plus figure} The signed allocation is not a monotone function
of radius, which is the point: at equal budget it \ph{describe the qualitative
structure --- where it spends relative to \Asens, and whether the departure is
phase-locked}. Reporting this matters because it is the part a deployable rule
has to imitate.

\begin{figure}[!tbp]
  \centering
  \ph{figures/fig\_oa\_schedules.pdf}
  \caption{Degree histories at $\beta=1$ on a representative orbit:
  \Fop, \Rint, \Asens\ and \Asign. The lower panel shows the accumulated
  displacement each produces.}
  \label{fig:schedules}
\end{figure}

\subsection{Certificate}

Table~\ref{tab:certificate} gives the certified gap. \ph{one paragraph}
\ph{State the median gap, the worst orbit, and therefore which of the two names
of Section~\ref{sec:algorithm-certificate} the paper uses.} Where the gap is
wide the reported ratios are conservative: the true optimum can only be better
than the schedule found, so the advantage over \Fop\ is understated and the
capture fraction of Section~\ref{sec:controller} overstated.

\subsection{Sensitivity of the benchmark to its own construction}

Three properties of the construction were audited and none moves a direction
reported above. \wpref{WP3, WP4}

\begin{itemize}
  \item \textbf{Start dependence.} Table~\ref{tab:starts}; the spread across
    \ph{n} starts is \ph{value}.
  \item \textbf{Accumulation grid.} Section~\ref{sec:setup-grids}; the schedule
    changes by \ph{value} between \ph{value}~s and \ph{value}~s.
  \item \textbf{Degree set.} \ph{one sentence} A coarser $\mathcal{N}$ can only
    make the benchmark worse, so the measured gap to the policies is
    conservative in that direction too.
\end{itemize}
